% Reciprocal proofs for the source-specific 1987 tutorial route.
\section{Klop 1987: abstract, residual, and strategy arguments}

\begin{annexitem}{annex:p12-klop1987-group-completion}{thm:p12-klop1987-group-completion}{The ten actual group rules: order, equational peaks, and every critical pair}
\begin{proof}
\paragraph{Termination with the stated comparison interface.}
Use the LPO of Definition~\ref{def:p12-ro-lpo} on the finite signature
with precedence \(I>_\Sigma .>_\Sigma e\), and ordinary left-to-right
status at the binary product. Theorem~\ref{thm:p12-ro-lpo-termination}
and its structural lemma give well-foundedness and preservation under
all contexts and substitutions. This is an editorial verification by
the available lexicographic-order interface, not a proof of the
report's separate marked-tree/Kruskal development.

Rules 1, 4, 8, 9, 10 and 12 have a proper subterm on the right.
Rules 2 and 11 decrease because their product head exceeds the
nullary \(e\). For rule 3, put \(s=(x.y).z\).
It exceeds \(x,y,z\) by the subterm property. It exceeds \(y.z\)
by the equal-head lexicographic clause, since \(x.y>y\) and
\(s>y,s>z\). It consequently exceeds \(x.(y.z)\), since
\(x.y>x\) and it dominates both target arguments.
For rule 15, put \(s=I(x.y)\).
It exceeds \(I(y)\) by the equal-unary-head clause, since
\(x.y>y\) and \(s>y\); likewise it exceeds \(I(x)\).
The precedence \(I>.\) and these two domination inequalities then
give \(I(x.y)>I(y).I(x)\). All ten rules therefore decrease.
The ordering's substitution and context properties apply these
comparisons to every actual rewrite step, proving SN.

\paragraph{Every retained equation follows from the three initial ones.}
The three initial rules 1, 2, 3 are retained. Derive the further rule
equations successively from their particular peaks, retaining the
source's temporary numbers even when their rules are later removed.
At each line both routes use only earlier derived equations, so its
two endpoints are equal in the original congruence:
\begin{align*}
4:\quad &(I(x).x).z \ \longrightarrow\ e.z\longrightarrow z,
 & &(I(x).x).z\longrightarrow I(x).(x.z);\\
5:\quad &I(I(y)).(I(y).y)\longrightarrow y,
 & &I(I(y)).(I(y).y)\longrightarrow I(I(y)).e;\\
6:\quad &I(e).(e.z)\longrightarrow z,
 & &I(e).(e.z)\longrightarrow I(e).z;\\
7:\quad &(I(I(y)).e).x\longrightarrow y.x,
 & &(I(I(y)).e).x\longrightarrow I(I(y)).(e.x)
     \longrightarrow I(I(y)).x;\\
8:\quad &I(I(y)).e\longrightarrow y,
 & &I(I(y)).e\longrightarrow y.e;\\
9:\quad &I(I(y)).e\longrightarrow y,
 & &I(I(y)).e\longrightarrow I(I(y));\\
10:\quad &I(e).e\longrightarrow e,
 & &I(e).e\longrightarrow I(e);\\
11:\quad &I(I(y)).I(y)\longrightarrow e,
 & &I(I(y)).I(y)\longrightarrow y.I(y);\\
12:\quad &(y.I(y)).x\longrightarrow e.x\longrightarrow x,
 & &(y.I(y)).x\longrightarrow y.(I(y).x);\\
13:\quad &(x.y).I(x.y)\longrightarrow e,
 & &(x.y).I(x.y)\longrightarrow x.(y.I(x.y));\\
14:\quad &I(x).[x.(y.I(x.y))]\longrightarrow y.I(x.y),
 & &I(x).[x.(y.I(x.y))]\longrightarrow I(x).e
     \longrightarrow I(x);\\
15:\quad &I(y).(y.I(x.y))\longrightarrow I(x.y),
 & &I(y).(y.I(x.y))\longrightarrow I(y).I(x).
\end{align*}
The congruence assertion does not depend on an unproved forward
orientation of any newly derived equality. Each equation is a
consequence of the earlier ones by transitivity and symmetry.
Thus all final rule equations belong to the initial theory; conversely
the initial three equations occur in the final rules. Their generated
congruences coincide. The cancellation of temporary rules 5, 6, 7,
13 and 14 does not discard their derived equalities from this argument.

\paragraph{Exhaustive nonvariable overlap enumeration.}
Rename the variables of the outer and inner rules apart.
For each of the ten outer left sides, visit its nonvariable positions;
unify that subterm with each of the ten inner left sides.
The outer left sides have respectively \(2,2,2,3,2,2,2,2,3,2\)
nonvariable positions, hence this gives exactly \(220\) candidates.
A variable-variable or variable-subterm overlap is not added to
this list; the ordinary critical-pair proof already treats that
separate substitution case, including repeated variables.

The following is the complete list of unifiable rule/position triples.
An entry \(j@p\) means inner rule \(j\) at outer position \(p\);
\(\epsilon\) is the root, and a string such as \(21\) means child 2
then child 1. The superscript \(\circ\) marks a literally identical
pair. These \(15\) cases join by zero steps. The other \(50\)
cases are listed with all joining steps below.
\begin{center}
\begin{tabular}{rl}
Outer rule & All unifiable inner-rule/position entries\\
\hline
1 & \(1@\epsilon^{\circ},\ 8@\epsilon^{\circ},\ 11@\epsilon,\ 12@\epsilon\)\\
2 & \(2@\epsilon^{\circ},\ 8@\epsilon,\ 9@1,\ 10@1,\ 15@1\)\\
3 & \(1@1,\ 2@1,\ 3@\epsilon^{\circ},\ 3@1,\ 4@1\)\\
 & \(8@\epsilon,\ 8@1,\ 11@\epsilon,\ 11@1,\ 12@\epsilon\)\\
 & \(12@1\)\\
4 & \(1@2,\ 2@2,\ 3@2,\ 4@\epsilon^{\circ},\ 4@2\)\\
 & \(8@2,\ 9@1,\ 10@1,\ 11@2,\ 12@2^{\circ}\)\\
 & \(15@1\)\\
8 & \(1@\epsilon^{\circ},\ 2@\epsilon,\ 3@\epsilon,\ 8@\epsilon^{\circ}\)\\
9 & \(9@\epsilon^{\circ},\ 9@1^{\circ},\ 10@1,\ 15@1\)\\
10 & \(10@\epsilon^{\circ}\)\\
11 & \(1@\epsilon,\ 3@\epsilon,\ 9@2,\ 10@2,\ 11@\epsilon^{\circ}\)\\
 & \(15@2\)\\
12 & \(1@\epsilon,\ 2@2,\ 3@\epsilon,\ 4@2^{\circ},\ 8@2\)\\
 & \(9@21,\ 10@21,\ 11@2,\ 12@\epsilon^{\circ},\ 12@2\)\\
 & \(15@21\)\\
15 & \(1@1,\ 2@1,\ 3@1,\ 4@1,\ 8@1\)\\
 & \(11@1,\ 12@1,\ 15@\epsilon^{\circ}\)\\
\end{tabular}
\end{center}

The finite unification calculation uses decomposition at equal function
heads, rejection of different heads or arities, and substitution of a
variable binding when its variable does not occur in the other term.
Each decomposition preserves exactly the solutions. Each permitted
binding factors every solution through that substitution; a proper
occurrence equation has no finite-tree solution. Binding eliminates a
variable, while decomposition between bindings removes function heads,
so the calculation terminates on the finite equation lists.
These steps prove that every accepted substitution is most general,
and every rejected candidate has no unifier; the table is therefore
an exhaustive overlap list, not a list of selected representative pairs.

For each remaining entry, \(u_j\) is its unified peak and
\(s_j,t_j\) are its two immediate reducts. All variables \(x_i\) are
fresh for that displayed pair. An arrow annotation \(r_k@p\) contracts
exactly source rule \(k\) at position \(p\).
The displayed common endpoint proves joinability whether or not
normal-form terminology is used. Each arrow is a direct instance of
one of the ten retained rules; context replacement at its annotated
position supplies its whole-term step.

\paragraph{Pair 1: outer 1, inner 11, position \(\epsilon\).}
\[
 u_{1}=(e\mathbin{.}I(e)).
\]
\begin{align*}
 s_{1}&=I(e),\\
 &\xrightarrow{r_{10}@\epsilon}e,\\
 t_{1}&=e.
\end{align*}

\paragraph{Pair 2: outer 1, inner 12, position \(\epsilon\).}
\[
 u_{2}=(e\mathbin{.}(I(e)\mathbin{.}x_{1})).
\]
\begin{align*}
 s_{2}&=(I(e)\mathbin{.}x_{1}),\\
 &\xrightarrow{r_{10}@1}(e\mathbin{.}x_{1}),\\
 &\xrightarrow{r_{1}@\epsilon}x_{1},\\
 t_{2}&=x_{1}.
\end{align*}

\paragraph{Pair 3: outer 2, inner 8, position \(\epsilon\).}
\[
 u_{3}=(I(e)\mathbin{.}e).
\]
\begin{align*}
 s_{3}&=e,\\
 t_{3}&=I(e),\\
 &\xrightarrow{r_{10}@\epsilon}e.
\end{align*}

\paragraph{Pair 4: outer 2, inner 9, position \(1\).}
\[
 u_{4}=(I(I(x_{1}))\mathbin{.}I(x_{1})).
\]
\begin{align*}
 s_{4}&=e,\\
 t_{4}&=(x_{1}\mathbin{.}I(x_{1})),\\
 &\xrightarrow{r_{11}@\epsilon}e.
\end{align*}

\paragraph{Pair 5: outer 2, inner 10, position \(1\).}
\[
 u_{5}=(I(e)\mathbin{.}e).
\]
\begin{align*}
 s_{5}&=e,\\
 t_{5}&=(e\mathbin{.}e),\\
 &\xrightarrow{r_{1}@\epsilon}e.
\end{align*}

\paragraph{Pair 6: outer 2, inner 15, position \(1\).}
\[
 u_{6}=(I((x_{1}\mathbin{.}x_{2}))\mathbin{.}(x_{1}\mathbin{.}x_{2})).
\]
\begin{align*}
 s_{6}&=e,\\
 t_{6}&=((I(x_{2})\mathbin{.}I(x_{1}))\mathbin{.}(x_{1}\mathbin{.}x_{2})),\\
 &\xrightarrow{r_{3}@\epsilon}(I(x_{2})\mathbin{.}(I(x_{1})\mathbin{.}(x_{1}\mathbin{.}x_{2}))),\\
 &\xrightarrow{r_{4}@2}(I(x_{2})\mathbin{.}x_{2}),\\
 &\xrightarrow{r_{2}@\epsilon}e.
\end{align*}

\paragraph{Pair 7: outer 3, inner 1, position \(1\).}
\[
 u_{7}=((e\mathbin{.}x_{1})\mathbin{.}x_{2}).
\]
\begin{align*}
 s_{7}&=(e\mathbin{.}(x_{1}\mathbin{.}x_{2})),\\
 &\xrightarrow{r_{1}@\epsilon}(x_{1}\mathbin{.}x_{2}),\\
 t_{7}&=(x_{1}\mathbin{.}x_{2}).
\end{align*}

\paragraph{Pair 8: outer 3, inner 2, position \(1\).}
\[
 u_{8}=((I(x_{1})\mathbin{.}x_{1})\mathbin{.}x_{2}).
\]
\begin{align*}
 s_{8}&=(I(x_{1})\mathbin{.}(x_{1}\mathbin{.}x_{2})),\\
 &\xrightarrow{r_{4}@\epsilon}x_{2},\\
 t_{8}&=(e\mathbin{.}x_{2}),\\
 &\xrightarrow{r_{1}@\epsilon}x_{2}.
\end{align*}

\paragraph{Pair 9: outer 3, inner 3, position \(1\).}
\[
 u_{9}=(((x_{1}\mathbin{.}x_{2})\mathbin{.}x_{3})\mathbin{.}x_{4}).
\]
\begin{align*}
 s_{9}&=((x_{1}\mathbin{.}x_{2})\mathbin{.}(x_{3}\mathbin{.}x_{4})),\\
 &\xrightarrow{r_{3}@\epsilon}(x_{1}\mathbin{.}(x_{2}\mathbin{.}(x_{3}\mathbin{.}x_{4}))),\\
 t_{9}&=((x_{1}\mathbin{.}(x_{2}\mathbin{.}x_{3}))\mathbin{.}x_{4}),\\
 &\xrightarrow{r_{3}@\epsilon}(x_{1}\mathbin{.}((x_{2}\mathbin{.}x_{3})\mathbin{.}x_{4})),\\
 &\xrightarrow{r_{3}@2}(x_{1}\mathbin{.}(x_{2}\mathbin{.}(x_{3}\mathbin{.}x_{4}))).
\end{align*}

\paragraph{Pair 10: outer 3, inner 4, position \(1\).}
\[
 u_{10}=((I(x_{1})\mathbin{.}(x_{1}\mathbin{.}x_{2}))\mathbin{.}x_{3}).
\]
\begin{align*}
 s_{10}&=(I(x_{1})\mathbin{.}((x_{1}\mathbin{.}x_{2})\mathbin{.}x_{3})),\\
 &\xrightarrow{r_{3}@2}(I(x_{1})\mathbin{.}(x_{1}\mathbin{.}(x_{2}\mathbin{.}x_{3}))),\\
 &\xrightarrow{r_{4}@\epsilon}(x_{2}\mathbin{.}x_{3}),\\
 t_{10}&=(x_{2}\mathbin{.}x_{3}).
\end{align*}

\paragraph{Pair 11: outer 3, inner 8, position \(\epsilon\).}
\[
 u_{11}=((x_{1}\mathbin{.}x_{2})\mathbin{.}e).
\]
\begin{align*}
 s_{11}&=(x_{1}\mathbin{.}(x_{2}\mathbin{.}e)),\\
 &\xrightarrow{r_{8}@2}(x_{1}\mathbin{.}x_{2}),\\
 t_{11}&=(x_{1}\mathbin{.}x_{2}).
\end{align*}

\paragraph{Pair 12: outer 3, inner 8, position \(1\).}
\[
 u_{12}=((x_{1}\mathbin{.}e)\mathbin{.}x_{2}).
\]
\begin{align*}
 s_{12}&=(x_{1}\mathbin{.}(e\mathbin{.}x_{2})),\\
 &\xrightarrow{r_{1}@2}(x_{1}\mathbin{.}x_{2}),\\
 t_{12}&=(x_{1}\mathbin{.}x_{2}).
\end{align*}

\paragraph{Pair 13: outer 3, inner 11, position \(\epsilon\).}
\[
 u_{13}=((x_{1}\mathbin{.}x_{2})\mathbin{.}I((x_{1}\mathbin{.}x_{2}))).
\]
\begin{align*}
 s_{13}&=(x_{1}\mathbin{.}(x_{2}\mathbin{.}I((x_{1}\mathbin{.}x_{2})))),\\
 &\xrightarrow{r_{15}@22}(x_{1}\mathbin{.}(x_{2}\mathbin{.}(I(x_{2})\mathbin{.}I(x_{1})))),\\
 &\xrightarrow{r_{12}@2}(x_{1}\mathbin{.}I(x_{1})),\\
 &\xrightarrow{r_{11}@\epsilon}e,\\
 t_{13}&=e.
\end{align*}

\paragraph{Pair 14: outer 3, inner 11, position \(1\).}
\[
 u_{14}=((x_{1}\mathbin{.}I(x_{1}))\mathbin{.}x_{2}).
\]
\begin{align*}
 s_{14}&=(x_{1}\mathbin{.}(I(x_{1})\mathbin{.}x_{2})),\\
 &\xrightarrow{r_{12}@\epsilon}x_{2},\\
 t_{14}&=(e\mathbin{.}x_{2}),\\
 &\xrightarrow{r_{1}@\epsilon}x_{2}.
\end{align*}

\paragraph{Pair 15: outer 3, inner 12, position \(\epsilon\).}
\[
 u_{15}=((x_{1}\mathbin{.}x_{2})\mathbin{.}(I((x_{1}\mathbin{.}x_{2}))\mathbin{.}x_{3})).
\]
\begin{align*}
 s_{15}&=(x_{1}\mathbin{.}(x_{2}\mathbin{.}(I((x_{1}\mathbin{.}x_{2}))\mathbin{.}x_{3}))),\\
 &\xrightarrow{r_{15}@221}(x_{1}\mathbin{.}(x_{2}\mathbin{.}((I(x_{2})\mathbin{.}I(x_{1}))\mathbin{.}x_{3}))),\\
 &\xrightarrow{r_{3}@22}(x_{1}\mathbin{.}(x_{2}\mathbin{.}(I(x_{2})\mathbin{.}(I(x_{1})\mathbin{.}x_{3})))),\\
 &\xrightarrow{r_{12}@2}(x_{1}\mathbin{.}(I(x_{1})\mathbin{.}x_{3})),\\
 &\xrightarrow{r_{12}@\epsilon}x_{3},\\
 t_{15}&=x_{3}.
\end{align*}

\paragraph{Pair 16: outer 3, inner 12, position \(1\).}
\[
 u_{16}=((x_{1}\mathbin{.}(I(x_{1})\mathbin{.}x_{2}))\mathbin{.}x_{3}).
\]
\begin{align*}
 s_{16}&=(x_{1}\mathbin{.}((I(x_{1})\mathbin{.}x_{2})\mathbin{.}x_{3})),\\
 &\xrightarrow{r_{3}@2}(x_{1}\mathbin{.}(I(x_{1})\mathbin{.}(x_{2}\mathbin{.}x_{3}))),\\
 &\xrightarrow{r_{12}@\epsilon}(x_{2}\mathbin{.}x_{3}),\\
 t_{16}&=(x_{2}\mathbin{.}x_{3}).
\end{align*}

\paragraph{Pair 17: outer 4, inner 1, position \(2\).}
\[
 u_{17}=(I(e)\mathbin{.}(e\mathbin{.}x_{1})).
\]
\begin{align*}
 s_{17}&=x_{1},\\
 t_{17}&=(I(e)\mathbin{.}x_{1}),\\
 &\xrightarrow{r_{10}@1}(e\mathbin{.}x_{1}),\\
 &\xrightarrow{r_{1}@\epsilon}x_{1}.
\end{align*}

\paragraph{Pair 18: outer 4, inner 2, position \(2\).}
\[
 u_{18}=(I(I(x_{1}))\mathbin{.}(I(x_{1})\mathbin{.}x_{1})).
\]
\begin{align*}
 s_{18}&=x_{1},\\
 t_{18}&=(I(I(x_{1}))\mathbin{.}e),\\
 &\xrightarrow{r_{8}@\epsilon}I(I(x_{1})),\\
 &\xrightarrow{r_{9}@\epsilon}x_{1}.
\end{align*}

\paragraph{Pair 19: outer 4, inner 3, position \(2\).}
\[
 u_{19}=(I((x_{1}\mathbin{.}x_{2}))\mathbin{.}((x_{1}\mathbin{.}x_{2})\mathbin{.}x_{3})).
\]
\begin{align*}
 s_{19}&=x_{3},\\
 t_{19}&=(I((x_{1}\mathbin{.}x_{2}))\mathbin{.}(x_{1}\mathbin{.}(x_{2}\mathbin{.}x_{3}))),\\
 &\xrightarrow{r_{15}@1}((I(x_{2})\mathbin{.}I(x_{1}))\mathbin{.}(x_{1}\mathbin{.}(x_{2}\mathbin{.}x_{3}))),\\
 &\xrightarrow{r_{3}@\epsilon}(I(x_{2})\mathbin{.}(I(x_{1})\mathbin{.}(x_{1}\mathbin{.}(x_{2}\mathbin{.}x_{3})))),\\
 &\xrightarrow{r_{4}@2}(I(x_{2})\mathbin{.}(x_{2}\mathbin{.}x_{3})),\\
 &\xrightarrow{r_{4}@\epsilon}x_{3}.
\end{align*}

\paragraph{Pair 20: outer 4, inner 4, position \(2\).}
\[
 u_{20}=(I(I(x_{1}))\mathbin{.}(I(x_{1})\mathbin{.}(x_{1}\mathbin{.}x_{2}))).
\]
\begin{align*}
 s_{20}&=(x_{1}\mathbin{.}x_{2}),\\
 t_{20}&=(I(I(x_{1}))\mathbin{.}x_{2}),\\
 &\xrightarrow{r_{9}@1}(x_{1}\mathbin{.}x_{2}).
\end{align*}

\paragraph{Pair 21: outer 4, inner 8, position \(2\).}
\[
 u_{21}=(I(x_{1})\mathbin{.}(x_{1}\mathbin{.}e)).
\]
\begin{align*}
 s_{21}&=e,\\
 t_{21}&=(I(x_{1})\mathbin{.}x_{1}),\\
 &\xrightarrow{r_{2}@\epsilon}e.
\end{align*}

\paragraph{Pair 22: outer 4, inner 9, position \(1\).}
\[
 u_{22}=(I(I(x_{1}))\mathbin{.}(I(x_{1})\mathbin{.}x_{2})).
\]
\begin{align*}
 s_{22}&=x_{2},\\
 t_{22}&=(x_{1}\mathbin{.}(I(x_{1})\mathbin{.}x_{2})),\\
 &\xrightarrow{r_{12}@\epsilon}x_{2}.
\end{align*}

\paragraph{Pair 23: outer 4, inner 10, position \(1\).}
\[
 u_{23}=(I(e)\mathbin{.}(e\mathbin{.}x_{1})).
\]
\begin{align*}
 s_{23}&=x_{1},\\
 t_{23}&=(e\mathbin{.}(e\mathbin{.}x_{1})),\\
 &\xrightarrow{r_{1}@\epsilon}(e\mathbin{.}x_{1}),\\
 &\xrightarrow{r_{1}@\epsilon}x_{1}.
\end{align*}

\paragraph{Pair 24: outer 4, inner 11, position \(2\).}
\[
 u_{24}=(I(x_{1})\mathbin{.}(x_{1}\mathbin{.}I(x_{1}))).
\]
\begin{align*}
 s_{24}&=I(x_{1}),\\
 t_{24}&=(I(x_{1})\mathbin{.}e),\\
 &\xrightarrow{r_{8}@\epsilon}I(x_{1}).
\end{align*}

\paragraph{Pair 25: outer 4, inner 15, position \(1\).}
\[
 u_{25}=(I((x_{1}\mathbin{.}x_{2}))\mathbin{.}((x_{1}\mathbin{.}x_{2})\mathbin{.}x_{3})).
\]
\begin{align*}
 s_{25}&=x_{3},\\
 t_{25}&=((I(x_{2})\mathbin{.}I(x_{1}))\mathbin{.}((x_{1}\mathbin{.}x_{2})\mathbin{.}x_{3})),\\
 &\xrightarrow{r_{3}@\epsilon}(I(x_{2})\mathbin{.}(I(x_{1})\mathbin{.}((x_{1}\mathbin{.}x_{2})\mathbin{.}x_{3}))),\\
 &\xrightarrow{r_{3}@22}(I(x_{2})\mathbin{.}(I(x_{1})\mathbin{.}(x_{1}\mathbin{.}(x_{2}\mathbin{.}x_{3})))),\\
 &\xrightarrow{r_{4}@2}(I(x_{2})\mathbin{.}(x_{2}\mathbin{.}x_{3})),\\
 &\xrightarrow{r_{4}@\epsilon}x_{3}.
\end{align*}

\paragraph{Pair 26: outer 8, inner 2, position \(\epsilon\).}
\[
 u_{26}=(I(e)\mathbin{.}e).
\]
\begin{align*}
 s_{26}&=I(e),\\
 &\xrightarrow{r_{10}@\epsilon}e,\\
 t_{26}&=e.
\end{align*}

\paragraph{Pair 27: outer 8, inner 3, position \(\epsilon\).}
\[
 u_{27}=((x_{1}\mathbin{.}x_{2})\mathbin{.}e).
\]
\begin{align*}
 s_{27}&=(x_{1}\mathbin{.}x_{2}),\\
 t_{27}&=(x_{1}\mathbin{.}(x_{2}\mathbin{.}e)),\\
 &\xrightarrow{r_{8}@2}(x_{1}\mathbin{.}x_{2}).
\end{align*}

\paragraph{Pair 28: outer 9, inner 10, position \(1\).}
\[
 u_{28}=I(I(e)).
\]
\begin{align*}
 s_{28}&=e,\\
 t_{28}&=I(e),\\
 &\xrightarrow{r_{10}@\epsilon}e.
\end{align*}

\paragraph{Pair 29: outer 9, inner 15, position \(1\).}
\[
 u_{29}=I(I((x_{1}\mathbin{.}x_{2}))).
\]
\begin{align*}
 s_{29}&=(x_{1}\mathbin{.}x_{2}),\\
 t_{29}&=I((I(x_{2})\mathbin{.}I(x_{1}))),\\
 &\xrightarrow{r_{15}@\epsilon}(I(I(x_{1}))\mathbin{.}I(I(x_{2}))),\\
 &\xrightarrow{r_{9}@1}(x_{1}\mathbin{.}I(I(x_{2}))),\\
 &\xrightarrow{r_{9}@2}(x_{1}\mathbin{.}x_{2}).
\end{align*}

\paragraph{Pair 30: outer 11, inner 1, position \(\epsilon\).}
\[
 u_{30}=(e\mathbin{.}I(e)).
\]
\begin{align*}
 s_{30}&=e,\\
 t_{30}&=I(e),\\
 &\xrightarrow{r_{10}@\epsilon}e.
\end{align*}

\paragraph{Pair 31: outer 11, inner 3, position \(\epsilon\).}
\[
 u_{31}=((x_{1}\mathbin{.}x_{2})\mathbin{.}I((x_{1}\mathbin{.}x_{2}))).
\]
\begin{align*}
 s_{31}&=e,\\
 t_{31}&=(x_{1}\mathbin{.}(x_{2}\mathbin{.}I((x_{1}\mathbin{.}x_{2})))),\\
 &\xrightarrow{r_{15}@22}(x_{1}\mathbin{.}(x_{2}\mathbin{.}(I(x_{2})\mathbin{.}I(x_{1})))),\\
 &\xrightarrow{r_{12}@2}(x_{1}\mathbin{.}I(x_{1})),\\
 &\xrightarrow{r_{11}@\epsilon}e.
\end{align*}

\paragraph{Pair 32: outer 11, inner 9, position \(2\).}
\[
 u_{32}=(I(x_{1})\mathbin{.}I(I(x_{1}))).
\]
\begin{align*}
 s_{32}&=e,\\
 t_{32}&=(I(x_{1})\mathbin{.}x_{1}),\\
 &\xrightarrow{r_{2}@\epsilon}e.
\end{align*}

\paragraph{Pair 33: outer 11, inner 10, position \(2\).}
\[
 u_{33}=(e\mathbin{.}I(e)).
\]
\begin{align*}
 s_{33}&=e,\\
 t_{33}&=(e\mathbin{.}e),\\
 &\xrightarrow{r_{1}@\epsilon}e.
\end{align*}

\paragraph{Pair 34: outer 11, inner 15, position \(2\).}
\[
 u_{34}=((x_{1}\mathbin{.}x_{2})\mathbin{.}I((x_{1}\mathbin{.}x_{2}))).
\]
\begin{align*}
 s_{34}&=e,\\
 t_{34}&=((x_{1}\mathbin{.}x_{2})\mathbin{.}(I(x_{2})\mathbin{.}I(x_{1}))),\\
 &\xrightarrow{r_{3}@\epsilon}(x_{1}\mathbin{.}(x_{2}\mathbin{.}(I(x_{2})\mathbin{.}I(x_{1})))),\\
 &\xrightarrow{r_{12}@2}(x_{1}\mathbin{.}I(x_{1})),\\
 &\xrightarrow{r_{11}@\epsilon}e.
\end{align*}

\paragraph{Pair 35: outer 12, inner 1, position \(\epsilon\).}
\[
 u_{35}=(e\mathbin{.}(I(e)\mathbin{.}x_{1})).
\]
\begin{align*}
 s_{35}&=x_{1},\\
 t_{35}&=(I(e)\mathbin{.}x_{1}),\\
 &\xrightarrow{r_{10}@1}(e\mathbin{.}x_{1}),\\
 &\xrightarrow{r_{1}@\epsilon}x_{1}.
\end{align*}

\paragraph{Pair 36: outer 12, inner 2, position \(2\).}
\[
 u_{36}=(x_{1}\mathbin{.}(I(x_{1})\mathbin{.}x_{1})).
\]
\begin{align*}
 s_{36}&=x_{1},\\
 t_{36}&=(x_{1}\mathbin{.}e),\\
 &\xrightarrow{r_{8}@\epsilon}x_{1}.
\end{align*}

\paragraph{Pair 37: outer 12, inner 3, position \(\epsilon\).}
\[
 u_{37}=((x_{1}\mathbin{.}x_{2})\mathbin{.}(I((x_{1}\mathbin{.}x_{2}))\mathbin{.}x_{3})).
\]
\begin{align*}
 s_{37}&=x_{3},\\
 t_{37}&=(x_{1}\mathbin{.}(x_{2}\mathbin{.}(I((x_{1}\mathbin{.}x_{2}))\mathbin{.}x_{3}))),\\
 &\xrightarrow{r_{15}@221}(x_{1}\mathbin{.}(x_{2}\mathbin{.}((I(x_{2})\mathbin{.}I(x_{1}))\mathbin{.}x_{3}))),\\
 &\xrightarrow{r_{3}@22}(x_{1}\mathbin{.}(x_{2}\mathbin{.}(I(x_{2})\mathbin{.}(I(x_{1})\mathbin{.}x_{3})))),\\
 &\xrightarrow{r_{12}@2}(x_{1}\mathbin{.}(I(x_{1})\mathbin{.}x_{3})),\\
 &\xrightarrow{r_{12}@\epsilon}x_{3}.
\end{align*}

\paragraph{Pair 38: outer 12, inner 8, position \(2\).}
\[
 u_{38}=(x_{1}\mathbin{.}(I(x_{1})\mathbin{.}e)).
\]
\begin{align*}
 s_{38}&=e,\\
 t_{38}&=(x_{1}\mathbin{.}I(x_{1})),\\
 &\xrightarrow{r_{11}@\epsilon}e.
\end{align*}

\paragraph{Pair 39: outer 12, inner 9, position \(21\).}
\[
 u_{39}=(I(x_{1})\mathbin{.}(I(I(x_{1}))\mathbin{.}x_{2})).
\]
\begin{align*}
 s_{39}&=x_{2},\\
 t_{39}&=(I(x_{1})\mathbin{.}(x_{1}\mathbin{.}x_{2})),\\
 &\xrightarrow{r_{4}@\epsilon}x_{2}.
\end{align*}

\paragraph{Pair 40: outer 12, inner 10, position \(21\).}
\[
 u_{40}=(e\mathbin{.}(I(e)\mathbin{.}x_{1})).
\]
\begin{align*}
 s_{40}&=x_{1},\\
 t_{40}&=(e\mathbin{.}(e\mathbin{.}x_{1})),\\
 &\xrightarrow{r_{1}@\epsilon}(e\mathbin{.}x_{1}),\\
 &\xrightarrow{r_{1}@\epsilon}x_{1}.
\end{align*}

\paragraph{Pair 41: outer 12, inner 11, position \(2\).}
\[
 u_{41}=(x_{1}\mathbin{.}(I(x_{1})\mathbin{.}I(I(x_{1})))).
\]
\begin{align*}
 s_{41}&=I(I(x_{1})),\\
 &\xrightarrow{r_{9}@\epsilon}x_{1},\\
 t_{41}&=(x_{1}\mathbin{.}e),\\
 &\xrightarrow{r_{8}@\epsilon}x_{1}.
\end{align*}

\paragraph{Pair 42: outer 12, inner 12, position \(2\).}
\[
 u_{42}=(x_{1}\mathbin{.}(I(x_{1})\mathbin{.}(I(I(x_{1}))\mathbin{.}x_{2}))).
\]
\begin{align*}
 s_{42}&=(I(I(x_{1}))\mathbin{.}x_{2}),\\
 &\xrightarrow{r_{9}@1}(x_{1}\mathbin{.}x_{2}),\\
 t_{42}&=(x_{1}\mathbin{.}x_{2}).
\end{align*}

\paragraph{Pair 43: outer 12, inner 15, position \(21\).}
\[
 u_{43}=((x_{1}\mathbin{.}x_{2})\mathbin{.}(I((x_{1}\mathbin{.}x_{2}))\mathbin{.}x_{3})).
\]
\begin{align*}
 s_{43}&=x_{3},\\
 t_{43}&=((x_{1}\mathbin{.}x_{2})\mathbin{.}((I(x_{2})\mathbin{.}I(x_{1}))\mathbin{.}x_{3})),\\
 &\xrightarrow{r_{3}@\epsilon}(x_{1}\mathbin{.}(x_{2}\mathbin{.}((I(x_{2})\mathbin{.}I(x_{1}))\mathbin{.}x_{3}))),\\
 &\xrightarrow{r_{3}@22}(x_{1}\mathbin{.}(x_{2}\mathbin{.}(I(x_{2})\mathbin{.}(I(x_{1})\mathbin{.}x_{3})))),\\
 &\xrightarrow{r_{12}@2}(x_{1}\mathbin{.}(I(x_{1})\mathbin{.}x_{3})),\\
 &\xrightarrow{r_{12}@\epsilon}x_{3}.
\end{align*}

\paragraph{Pair 44: outer 15, inner 1, position \(1\).}
\[
 u_{44}=I((e\mathbin{.}x_{1})).
\]
\begin{align*}
 s_{44}&=(I(x_{1})\mathbin{.}I(e)),\\
 &\xrightarrow{r_{10}@2}(I(x_{1})\mathbin{.}e),\\
 &\xrightarrow{r_{8}@\epsilon}I(x_{1}),\\
 t_{44}&=I(x_{1}).
\end{align*}

\paragraph{Pair 45: outer 15, inner 2, position \(1\).}
\[
 u_{45}=I((I(x_{1})\mathbin{.}x_{1})).
\]
\begin{align*}
 s_{45}&=(I(x_{1})\mathbin{.}I(I(x_{1}))),\\
 &\xrightarrow{r_{11}@\epsilon}e,\\
 t_{45}&=I(e),\\
 &\xrightarrow{r_{10}@\epsilon}e.
\end{align*}

\paragraph{Pair 46: outer 15, inner 3, position \(1\).}
\[
 u_{46}=I(((x_{1}\mathbin{.}x_{2})\mathbin{.}x_{3})).
\]
\begin{align*}
 s_{46}&=(I(x_{3})\mathbin{.}I((x_{1}\mathbin{.}x_{2}))),\\
 &\xrightarrow{r_{15}@2}(I(x_{3})\mathbin{.}(I(x_{2})\mathbin{.}I(x_{1}))),\\
 t_{46}&=I((x_{1}\mathbin{.}(x_{2}\mathbin{.}x_{3}))),\\
 &\xrightarrow{r_{15}@\epsilon}(I((x_{2}\mathbin{.}x_{3}))\mathbin{.}I(x_{1})),\\
 &\xrightarrow{r_{15}@1}((I(x_{3})\mathbin{.}I(x_{2}))\mathbin{.}I(x_{1})),\\
 &\xrightarrow{r_{3}@\epsilon}(I(x_{3})\mathbin{.}(I(x_{2})\mathbin{.}I(x_{1}))).
\end{align*}

\paragraph{Pair 47: outer 15, inner 4, position \(1\).}
\[
 u_{47}=I((I(x_{1})\mathbin{.}(x_{1}\mathbin{.}x_{2}))).
\]
\begin{align*}
 s_{47}&=(I((x_{1}\mathbin{.}x_{2}))\mathbin{.}I(I(x_{1}))),\\
 &\xrightarrow{r_{15}@1}((I(x_{2})\mathbin{.}I(x_{1}))\mathbin{.}I(I(x_{1}))),\\
 &\xrightarrow{r_{3}@\epsilon}(I(x_{2})\mathbin{.}(I(x_{1})\mathbin{.}I(I(x_{1})))),\\
 &\xrightarrow{r_{11}@2}(I(x_{2})\mathbin{.}e),\\
 &\xrightarrow{r_{8}@\epsilon}I(x_{2}),\\
 t_{47}&=I(x_{2}).
\end{align*}

\paragraph{Pair 48: outer 15, inner 8, position \(1\).}
\[
 u_{48}=I((x_{1}\mathbin{.}e)).
\]
\begin{align*}
 s_{48}&=(I(e)\mathbin{.}I(x_{1})),\\
 &\xrightarrow{r_{10}@1}(e\mathbin{.}I(x_{1})),\\
 &\xrightarrow{r_{1}@\epsilon}I(x_{1}),\\
 t_{48}&=I(x_{1}).
\end{align*}

\paragraph{Pair 49: outer 15, inner 11, position \(1\).}
\[
 u_{49}=I((x_{1}\mathbin{.}I(x_{1}))).
\]
\begin{align*}
 s_{49}&=(I(I(x_{1}))\mathbin{.}I(x_{1})),\\
 &\xrightarrow{r_{2}@\epsilon}e,\\
 t_{49}&=I(e),\\
 &\xrightarrow{r_{10}@\epsilon}e.
\end{align*}

\paragraph{Pair 50: outer 15, inner 12, position \(1\).}
\[
 u_{50}=I((x_{1}\mathbin{.}(I(x_{1})\mathbin{.}x_{2}))).
\]
\begin{align*}
 s_{50}&=(I((I(x_{1})\mathbin{.}x_{2}))\mathbin{.}I(x_{1})),\\
 &\xrightarrow{r_{15}@1}((I(x_{2})\mathbin{.}I(I(x_{1})))\mathbin{.}I(x_{1})),\\
 &\xrightarrow{r_{3}@\epsilon}(I(x_{2})\mathbin{.}(I(I(x_{1}))\mathbin{.}I(x_{1}))),\\
 &\xrightarrow{r_{2}@2}(I(x_{2})\mathbin{.}e),\\
 &\xrightarrow{r_{8}@\epsilon}I(x_{2}),\\
 t_{50}&=I(x_{2}).
\end{align*}

All endpoints in each displayed pair coincide. Any instance of its
MGU further substitutes these same joining paths, by closure of
ordinary rewriting under substitution. Thus every critical pair
joins. The full ordinary criterion
(Theorem~\ref{thm:p12-critical-pair}) gives WCR, and the already
proved SN and the abstract Newman joining argument give CR.
All terms therefore have unique normal forms. Equality in the
initial congruence is finite conversion in the final rules, and CR
identifies convertible normal forms; conversely equal normal forms
give a conversion. Effective matching of the finite rules and
termination compute these normal forms, proving the original
ground word-problem claim as well as the displayed all-term
equational equivalence.
\end{proof}
\end{annexitem}


\begin{annexitem}{annex:p12-klop1987-ars-commutation}{thm:p12-klop1987-ars-commutation}{The finite strip inductions, with their short arms retained}
\begin{proof}
For Exercise (3), first take a single \(a\to_1 b\) and a finite
\(a\to_2^*c\). Induct on the latter path's length to obtain
\(b\to_2^*d,c\to_1^=d\). For length zero take \(d=b\).
Otherwise write \(a\to_2u\to_2^*c\). The assumed one-step
peak condition gives \(b\to_2^*v,u\to_1^=v\).
If \(u=v\), take \(d=c\), following \(v=u\to_2^*c\)
from \(b\). If \(u\to_1v\), apply the shorter-path induction
at \(u\) to its \(1\)-step to \(v\) and its \(2\)-path
to \(c\). This gives \(v\to_2^*d,c\to_1^=d\).
Concatenating the \(2\)-paths proves the assertion.
Now induct on the length of \(a\to_1^*b\), allowing the other
\(2\)-path any finite length. For the first \(1\)-step
\(a\to_1u\), the preceding strip supplies
\(u\to_2^*v,c\to_1^=v\). Apply the induction hypothesis
at \(u\) to its remaining \(1\)-path to \(b\) and its
\(2\)-path to \(v\). It joins \(b\) by \(2^*\) and
\(v\) by \(1^*\). Concatenation gives commutation.

For (4), start instead with a finite \(a\to_1^*b\) and one
\(a\to_2c\), and induct on the first length. The local condition
and the induction at the next \(1\)-vertex give
\(b\to_2d,c\to_1^*d\); the \(2\)-arm remains a single
step. The zero-length case is the original \(2\)-step.
Extend to an arbitrary \(2\)-path by induction on its length:
move its first step past the finite \(1\)-path using this strip,
then apply the induction to the shorter remaining \(2\)-path and
the resulting \(1\)-path. This gives the claimed commuting closures.
No bound on a newly formed long \(1\)-path was used as the
decreasing induction parameter.

For (10), take the argument for (3) with both relations equal to
\(\to\). Its hypothesis is exactly strong confluence, and its
conclusion is CR. Subcommutativity has at most one step on each arm,
so implies this hypothesis and CR. Thus in (1), \(\to_1\) is CR.
Equality \(\to_1^*=\to_2^*\) identifies all peaks and joins
at the finite-reduction level, so transports CR to \(\to_2\).

For (2), fix a colour \(\alpha\). Induct on the length of a union
path to show that \(\to_\alpha^*\) commutes with the union closure.
Given \(a\to_\alpha^*b\) and
\(a\to_\beta u\to^*c\), pairwise commutation gives
\(b\to_\beta^*v,u\to_\alpha^*v\). Apply the shorter-union-path
induction at \(u\) to obtain \(v\to^*d,c\to_\alpha^*d\),
and concatenate from \(b\). The zero-length case is immediate.
Finally induct on the length of one arbitrary union path in a peak:
move its first coloured step past the other path using the just-proved
strip, then join the remaining shorter path. This gives confluence
of the union. Only finitely many colours are used in each actual path;
the proof does not assume a finite index set.
\end{proof}
\end{annexitem}

\begin{annexitem}{annex:p12-klop1987-ars-request}{thm:p12-klop1987-ars-request}{A diamond relation made from the two full closure blocks}
\begin{proof}
Put \(P=\to_1^*\circ\to_2^*\), with composition in the
source's traversal order. It is reflexive, contains each original
one-step relation, and is contained in their union closure. Hence
\(P^*\) is exactly that union closure.

First suppose \(\to_1\) and \(\to_2\) are CR. In a \(P\)-peak
write \(a\to_1^*p\to_2^*b\) and
\(a\to_1^*q\to_2^*c\). Confluence of \(1\) gives
\(p\to_1^*r,q\to_1^*r\). Apply full requesting to the
peak from \(p\) ending at \(r,b\): there are
\(r\to_2^*d\) and \(b\to_1^*e\to_2^*d\).
The analogous peak from \(q\) gives
\(r\to_2^*d'\), \(c\to_1^*e'\to_2^*d'\).
Confluence of \(2\) joins \(d,d'\) at \(z\).
Thus \(bPz\) and \(cPz\), so \(P\) has the diamond property.
The finite grid induction proves confluence of \(P^*\), giving (5).

For (7), replace CR of \(1\) by the given \(1^*\)-peak condition.
It yields \(pP r,qP r\), say
\(p\to_1^*p'\to_2^*r\) and
\(q\to_1^*q'\to_2^*r\). Requesting at \(p\) gives
\(p'\to_2^*u\), \(b\to_1^*e\to_2^*u\).
Requesting at \(q\) gives
\(q'\to_2^*u'\), \(c\to_1^*e'\to_2^*u'\).
Use CR of \(2\) to join \(r,u\), both descendants of \(p'\),
at \(v\); join \(r,u'\), both descendants of \(q'\), at
\(v'\); then join \(v,v'\), both descendants of \(r\),
at \(z\). Concatenation again gives \(bPz,cPz\).
The same diamond and closure argument proves union CR. Each use of
requesting has two finite top arms and finite \(2\)-joining arms;
none is replaced by the local premise of (6).
\end{proof}
\end{annexitem}

\begin{annexitem}{annex:p12-klop1987-ars-refinement}{thm:p12-klop1987-ars-refinement}{Compatibility for every finite refining path}
\begin{proof}
Refinement implies \(\to_1^*\subseteq\to_2^*\), by replacing
each \(1\)-step by its finite \(2\)-path. If compatibility holds
and \(a\to_2b\to_1^*c\), this inclusion gives
\(a\to_2^*c\). Compatibility at that pair gives the common
\(1\)-descendant of \(a,c\) required in (8).

Conversely assume that mixed condition. Given a finite
\(a=a_0\to_2a_1\to_2\cdots\to_2a_n=b\) and a further
\(b\to_1^*c\), propagate the endpoint backwards. At the last
step the condition supplies \(c\to_1^*d_{n-1}\) and
\(a_{n-1}\to_1^*d_{n-1}\). At the preceding step apply it
with that latter path, obtaining a further \(1\)-descendant
\(d_{n-2}\) of \(d_{n-1}\) and of \(a_{n-2}\).
Continuing through the finite list gives a common \(1\)-descendant
of \(a,c\). With \(c=b\) and the empty further path this
is compatibility. For \(n=0\), take the given \(c\).

If \(1\) is CR and \(a\to_2^*b,a\to_2^*c\), compatibility
joins \(a,b\) by \(1\) at \(p\) and \(a,c\) by \(1\)
at \(q\). CR of \(1\) joins \(p,q\) by \(1\) at \(r\).
Refinement turns the resulting \(b\to_1^*r,c\to_1^*r\)
into a \(2\)-join, proving CR of \(2\).
For the converse, a \(1\)-peak is a \(2\)-peak by refinement;
join its endpoints \(b,c\) by \(2\) at \(d\). Compatibility
on \(b\to_2^*d\) gives \(b\to_1^*p,d\to_1^*p\).
The latter path refines to \(2\), so \(c\to_2^*p\).
Compatibility at this pair gives \(c\to_1^*q,p\to_1^*q\).
Then \(b\to_1^*p\to_1^*q\) and \(c\to_1^*q\), proving
CR of \(1\) without presupposing it in an intermediate join.

For the one-step compatibility warning, take
\(\to_1=\{a\to d,b\to c,b\to d\}\) and
\(\to_2=\to_1\cup\{a\to b\}\) on four distinct elements.
This is a refinement. Every \(2\)-step has a common \(1\)-descendant
of its endpoints: an original \(1\)-step joins at its target, and
\(a\to_2b\) joins at \(d\). But \(a\to_2b\to_2c\)
has no such join, since \(a\)'s only proper \(1\)-descendant
is \(d\) and \(c,d\) are distinct \(1\)-normal forms.
The mixed criterion also fails at \(a\to_2b\to_1c\).
\end{proof}
\end{annexitem}

\begin{annexitem}{annex:p12-klop1987-ars-splitting}{thm:p12-klop1987-ars-splitting}{Finite-colour partition into strongly commuting and terminating parts}
\begin{proof}
Fix an arbitrary finite union peak and let \(J\) be its finite set
of colours. Weak commutation of two colours joins by paths using
only those same colours. Thus all the constructions below stay inside
the subfamily indexed by \(J\).
Partition \(J\) into the globally splitting colours \(B\) and
the nonsplitting colours \(G\). Let \(S\) and \(T\) be their
respective unions. An infinite \(S\)-reduction would use some colour
of the finite set \(B\) infinitely often, contradicting that colour's
relative termination in the whole union. Hence \(S\) is SN.
Every one-step \(S\)-peak is weakly joined by its two colours, so
\(S\) is WCR. The well-founded joining induction now gives CR:
join the first steps, use the induction assertion at each strict
descendant to join its remaining endpoint with that first join, then
use the assertion at the first join to join the two resulting terms.
All these induction sources are strict descendants of the original
source; zero-length arms are immediate. This is the abstract Newman
argument and requires no term structure.

For a nonsplitting colour \(\gamma\), weak commutation supplies
at least one join of any mixed peak having that colour on an arm.
Nonsplitting says a join can be chosen with at most one step on the
other-coloured arm: otherwise that peak would witness splitting.
For a \(T\)-peak, select the colour on the second arm this way.
Its join has a finite \(T\)-arm from the first endpoint and at
most one \(T\)-step from the second. Thus \(T\) is strongly
confluent and CR by
Theorem~\ref{thm:p12-klop1987-ars-commutation}.

For a mixed peak \(a\to_Sb,a\to_Tc\), use nonsplitting of
the colour on the \(T\)-arm. It gives
\(b\to_T^*d,c\to_S^=d\), exactly the short-arm hypothesis
of Exercise (3) for the two union relations \(S,T\).
Their closures commute by the already proved strip argument.
They are individually CR, so their self-pairs also commute. The
pairwise-commuting union theorem proves CR of \(S\cup T\).
Our chosen peak consequently joins using only colours of \(J\).
Since the peak was arbitrary, the full, possibly infinitely indexed
union is CR. No argument asserted that an infinite union of relatively
terminating colours is itself terminating; finiteness of \(J\)
is used exactly at the pigeonhole step.
\end{proof}
\end{annexitem}

\begin{annexitem}{annex:p12-klop1987-ars-decreasing}{thm:p12-klop1987-ars-decreasing}{Well-founded joining along the entire decreasing conversion}
\begin{proof}
CR supplies a join \(b\to^*z,c\to^*z\) for every one-step
peak from \(a\). Follow the first path and reverse the second to
form a conversion from \(b\) to \(c\). Every term on it is
below \(a\): each endpoint is below \(a\), and a further
reduction stays below it by transitivity of \(>\).

Conversely perform well-founded induction on \(a\), proving that
any two finite reductions from \(a\) join. In a one-step peak
\(a\to u,a\to v\), its stipulated conversion has all vertices
strictly below \(a\). Join its successive vertices inductively along
the finite conversion: maintain a common descendant of \(u\) and
the current vertex. For a forward next step, the current vertex's
induction assertion joins that common descendant with the next vertex.
For a backward next step, the next vertex already reaches the current
one and its common descendant. Concatenation updates the join. At the
last vertex this gives a common descendant \(w\) of \(u,v\).

For arbitrary nonempty reductions from \(a\), take their first
steps \(u,v\) and this join \(w\). Apply the induction assertion
at \(u\) to join the first endpoint with \(w\), and at
\(v\) to join the second endpoint with \(w\). Finally apply
it at \(w\) to join those two resulting descendants. All three
sources are below \(a\), because each is a descendant of \(u\)
or \(v\) and the order contains all steps. If an original reduction
is empty its other endpoint is already a join. This proves the induction
assertion and CR. The proof uses every conversion vertex's strict bound,
not merely a bound on \(u,v\).
\end{proof}
\end{annexitem}

\begin{annexitem}{annex:p12-klop1987-ars-boundaries}{thm:p12-klop1987-ars-boundaries}{Successor closure and lifting through a retraction}
\begin{proof}
Every finite reduction starting in a successor-closed subset stays in
it. A member is normal in the restricted system exactly when it is
normal in the larger one. Thus normalizing paths, common reducts of
finite or one-step peaks, and the at-most-one-step joins stay in the
subset, proving inheritance of WN, CR, WCR and subcommutativity.
UN and NF follow using that same normal-form equivalence and inclusion
of every restricted conversion in the larger conversion relation.
Every restricted reduction sequence has the larger system's common
descendant under Ind; it is below its initial element, so lies in the
closed subset as well. SN, the increasing rank and the finite bound
on immediate successors restrict directly. This proves all ten
inheritance claims without treating an arbitrary induced subgraph as
a successor-closed one.

CR joins any convertible pair; irreducibility therefore prevents two
distinct normal forms from being convertible. Such a pair witnesses
consistency. If \(\mathcal A\) is consistent and its extension
is conservative, a witnessing nonconvertible pair from \(A\)
remains nonconvertible in \(B\), proving consistency of the extension.

Suppose instead the extension \(\mathcal B\) is confluent and
\(a,a'\in A\) are convertible in it. CR gives a common descendant
\(d\). Successor closure of \(A\) places every vertex of the
two joining paths, and \(d\), in \(A\); restriction of the
relation makes them \(\mathcal A\)-paths. Thus \(a,a'\) were
already convertible in \(\mathcal A\). Conversely each \(A\)-step
is a \(B\)-step, proving conservativity in both directions.

For (17), an infinite \(\alpha\)-sequence
\(a_0\to_\alpha a_1\to_\alpha\cdots\) lifts from
\(b_0=\iota(a_0)\), since \(\kappa(b_0)=a_0\).
At stage \(i\), the stipulated hypothesis for \emph{every}
\(b_i\) over \(a_i\) supplies an actual step
\(b_i\to_\beta b_{i+1}\) with \(\kappa(b_{i+1})=a_{i+1}\).
Continuing recursively gives an infinite \(\beta\)-sequence,
contrary to SN of \(\mathcal B\). Identity-valued genuine steps
still count; a lift allowed to be empty could not produce this contradiction.
\end{proof}
\end{annexitem}

\begin{annexitem}{annex:p12-klop1987-ars-normalization}{thm:p12-klop1987-ars-normalization}{Two distinct terminating ranks and the inverse-branching bound}
\begin{proof}
For the strict nontrivial diamond, prove by induction on \(n\) that
if \(a\) has a normalizing reduction of length \(n\), every
reduction from it has length at most \(n\). Length zero is a normal
form. Choose the first step \(a\to c\) of the given path.
If another first step has the same endpoint, the induction at \(c\)
gives the bound. If it has distinct endpoint \(b\), the strict
diamond gives \(c\to d,b\to d\). The induction at \(c\)
implies that \(d\) reaches a normal form in at most \(n-2\)
steps: continue a reduction from \(d\) maximally; a longer one
would extend a reduction from \(c\) beyond its bound, and a maximal
bounded reduction ends in normal form. Thus \(b\) has a normalizing
path of at most \(n-1\) steps, and the induction assertion at
\(b\) bounds every continuation by \(n-1\). For \(n=1\)
the distinct case is impossible, since \(c\) is normal but would
have a step to \(d\). This proves the bound, excluding infinite
reductions. WN supplies such an \(n\) for every element.

Now assume WCR and Inc, and fix a normal form \(z\) below \(a\).
Let \(h\) be the given increasing natural-number rank. Induct on
\(h(z)-h(a)\) to prove that every descendant of \(a\) reaches
\(z\). At zero, any positive reduction to \(z\) would increase
rank, so \(a=z\), and there is nothing to prove. Otherwise choose
the first step \(a\to c\) of a normalizing path to \(z\).
The rank difference at \(c\) is smaller; its induction assertion
says every descendant of \(c\) reaches \(z\). For any other
\(a\to b\), WCR joins \(b,c\) at \(d\).
The assertion at \(c\) gives \(d\to^*z\), so
\(b\to^*d\to^*z\). The rank difference at \(b\) is now
nonnegative and strictly smaller than at \(a\), so its induction
assertion controls all its descendants as well. These first-step cases
prove the claim. Therefore every descendant has rank at most \(h(z)\),
and an infinite increasing-rank reduction from \(a\) is impossible.
WN supplies \(z\) for every \(a\), proving Theorem 1.6(vi).

Inverting a step makes an Inc rank strictly decrease in \(\mathbb N\),
so Inc implies \(\mathrm{SN}^{-1}\). Conversely assume inverse
finite branching and inverse SN. At each \(a\), its finitely
branching predecessor-path tree has bounded depth: otherwise choose,
at each level, a child whose depths are unbounded to obtain an infinite
inverse reduction. Let \(h(a)\) be the maximum length of an inverse
path beginning at \(a\), a finite natural number. If \(a\to b\),
the inverse step \(b\to^{-1}a\) prefixes every inverse path
from \(a\); hence \(h(b)\geq h(a)+1\). This constructs Inc.
Forward finite branching is not the hypothesis used here.
\end{proof}
\end{annexitem}

\begin{annexitem}{annex:p12-klop1987-ars-countercases}{thm:p12-klop1987-ars-countercases}{Three finite failures with every premise checked}
\begin{proof}
For (6), use distinct \(a,b,c,d\),
\(\to_1=\{b\to a,c\to b\}\), and
\(\to_2=\{b\to c,c\to d\}\).
The second relation is terminating and confluent: it is the deterministic
chain \(b\to c\to d\), with \(a\) isolated.
There are exactly two mixed one-step peaks. At \(b\), their
endpoints are \(a,c\); choose joining endpoint \(a\), using
\(a\to_2^*a\) and \(c\to_1b\to_1a\to_2^*a\).
At \(c\), their endpoints are \(b,d\); choose joining endpoint
\(d\), using \(b\to_2c\to_2d\) and the empty
\(d\to_1^*d\to_2^*d\). Thus the displayed local premise
holds at every source. But \(b\to_1a\) and
\(b\to_2c\to_2d\) form a full peak whose endpoints
\(a,d\) are normal in both relations. They admit no requested
join. This checks the nonimplication without adding any missing
termination or confluence premise; even \(1\) is terminating and CR.

For (16), take one element \(x\), an \(\alpha\)-self-loop,
and no \(\beta\)-steps. The two-step quasi-commutation premise
has no instance and is true. \(\beta\) is SN, while the loop
gives an infinite \(\alpha\)-reduction. Hence the claimed SN
equivalence does not follow from that premise.

For the weaker version of (15), take \(x\to x,x\to y\),
with \(y\) normal. Both elements have normal forms. Every peak
joins at \(y\) by at most one step on each arm: use \(x\to y\)
from an \(x\)-endpoint and zero steps from a \(y\)-endpoint.
The repeated self-loop is infinite. Thus WN and
\(\mathrm{WCR}^{\leq1}\) do not imply SN. The distinct-endpoint
peak cannot have a strict step from \(y\), so the stricter positive
hypothesis is correctly excluded.
\end{proof}
\end{annexitem}

\begin{annexitem}{annex:p12-klop1987-leftmost-needed}{thm:p12-klop1987-leftmost-needed}{Persistence of the leftmost occurrence until its contraction}
\begin{proof}
Let \(u\) be the leftmost occurrence in the non-normal input, and
take any finite reduction to a normal form. Suppose no residual of
\(u\) is contracted in that reduction. The left-normal position
fact proved in Annex~\ref{annex:p12-klop1987-strategy-results} says
that every step at a different, non-leftmost redex preserves the unique
leftmost residual of \(u\). Apply that fact inductively through the
given finite reduction. At every stage \(u\) has a unique surviving
leftmost residual: no other redex can be above it and no step to its
right can erase, copy, or cover it. Its presence at the purported
normal-form endpoint contradicts irreducibility. Thus some residual
is contracted in every such reduction, which is exactly neededness.
Finite position search returns the leftmost redex when matching is
effective. Nothing in the argument decides whether its input has a
normal form, and nothing transports this selector beyond left-normality.
\end{proof}
\end{annexitem}

\begin{annexitem}{annex:p12-klop1987-outermost-projection}{thm:p12-klop1987-outermost-projection}{The ordinary disjoint-residual staircase and its nonempty tail}
\begin{proof}
Throughout, residuals of one designated original redex form a finite
disjoint set at every finite stage. This is the first-order regular
property; it is not assumed for calculi with bound variables.
We use the following two position facts and give their reasons.

\paragraph{A non-outermost step cannot create a new outermost redex.}
Let the contracted redex lie properly inside an existing redex \(p\).
Regularity puts the contraction inside a variable instance of \(p\),
so the left-linear match of \(p\) persists. Any redex created inside
it remains covered. A newly created redex above \(p\) would either
overlap \(p\)'s operator pattern, which is forbidden, or contain
\(p\) inside a variable instance. In the latter case its operator
pattern lies outside \(p\) and cannot be newly matched by a change
inside \(p\); left-linearity imposes no additional equality test.
Disjoint regions are unchanged. Thus every outermost redex after the
step has an outermost ancestor before it. Pulling back through finitely
many non-outermost steps yields an outermost trace at every intermediate
term.

\paragraph{Outermost traces exchange across an elementary commuting cell.}
Take two steps from \(a\), and complete their residual diamond.
Suppose a surviving redex trace is outermost along one route from
\(a\) to the common endpoint. Its final occurrence has a unique
residual ancestry on the other route: disjoint replacements preserve
positions; in a variable instance the chosen final copy determines the
original instance and copy; cancellation of its own contraction cannot
occur on a route where it survives at the endpoint. This checks the
disjoint, identical-root and variable-nested diamond cases.
If its ancestor on the other route were not outermost, it would be
covered by a redex newly created by that route's first contraction.
The old traced occurrence was outermost at \(a\), so the covering
redex did not exist there. The rest of the elementary route contracts
only residuals of the other original step, not this newly created
covering redex. Such a residual also cannot be above the trace: its
original ancestor would then have covered the original outermost
occurrence. The cover therefore persists to the common endpoint,
contradicting outermostness there. The exchanged trace is outermost.
Applying this check one cell at a time gives the exchange for finite
commuting strips. Reordering a finite disjoint development is justified
by those same cells, so the traced endpoint is unchanged.

Now let \(\rho:a=a_0\to a_1\to\cdots\) be outermost-fair,
and let its projection through a fixed redex step be
\(\rho':b=b_0\to^*b_1\to^*\cdots\). Write
\(D_i:a_i\to^*b_i\) for the development of the designated residual
set \(U_i\). Suppose \(\rho'\) had an infinite outermost trace.
Choose a boundary \(b_k\) on this trace; ignoring a finite prefix
does not affect the contradiction.

Reorder \(D_k\) by first contracting designated residuals which are
currently outermost, stopping when none remains, then contracting the
remaining designated residuals. Write the result as
\(a_k\xrightarrow{P_k}c_k\xrightarrow{Q_k}b_k\).
Both parts are finite developments with the same endpoint as \(D_k\).
Every step of \(Q_k\) is non-outermost, by the first position fact.
The outermost trace at \(b_k\) therefore pulls back to an outermost
trace from some occurrence at \(c_k\).

Project \(P_k\) past \(a_k\to a_{k+1}\), obtaining a development
from \(a_{k+1}\) to \(d_{k+1}\). Exchange the outermost trace
through the finite commuting strip. It yields an outermost trace through
\(d_{k+1}\) to \(b_{k+1}\). Reorder this projected development
again into an outermost prefix \(P_{k+1}\), ending at \(c_{k+1}\),
and a non-outermost remainder ending at \(d_{k+1}\). Pull back through
that remainder. This gives an outermost occurrence at \(c_{k+1}\)
and continues the trace to \(b_{k+1}\). Iterate this explicit procedure
over successive rows of the projection strip.

The prefix \(P_{i+1}\) develops only residuals of the labels developed
by \(P_i\); no new origin enters this sequence of prefixes. Moreover
each label in \(P_i\) was outermost already in \(a_i\).
Indeed those labels are disjoint. Contracting another disjoint label
cannot remove a pre-existing redex covering this label; it can only
leave its covering status unchanged or create an additional cover.
Thus a designated label selected as outermost later in the prefix could
not have been covered at that row's start.

If these prefixes stayed nonempty at arbitrarily late rows, their finite
initial label set and finite copying at each step would give, by the
finitely branching tree argument, an infinite label branch. Every
member of that branch is outermost in the corresponding original
\(a_i\), and every transition is a residual transition of
\(a_i\to a_{i+1}\). This is an infinite outermost trace in
\(\rho\), contradicting fairness. Consequently some \(P_l\)
is empty. All subsequent prefixes are empty as well, and
\(c_i=a_i\) from that row onward. The traced occurrences just
constructed through the \(c_i\) now give an infinite outermost trace
along the original steps of \(\rho\), the same contradiction.
This proves outermost fairness of the projection.

Next suppose \(\rho\) has infinitely many outermost contractions
but its projection is finite. Beyond some row all projected segments
are empty. A contraction disjoint from the designated set, or above
some of it, would project to a genuine step, so cannot occur in this
tail. A contraction strictly inside one of its designated residuals
may project to an empty segment by erasure, but is not outermost; it
does not increase the number of designated residuals, since those
residuals are disjoint and it contains none of them. Every outermost
step in this tail must therefore contract a designated residual itself,
decreasing their finite number by one. Infinitely many such decreases
are impossible. Equivalently, once that number is zero all subsequent
segments are the original genuine steps. The projection is infinite.

Finally, an infinite outermost-fair reduction cannot have only finitely
many outermost contractions. After the last such contraction choose an
outermost redex in the current non-normal term. All subsequent steps
are inside variable instances of that redex or disjoint from it; its
pattern persists. The first position fact prevents any newly created
outermost redex from covering it in a non-outermost step. It therefore
persists as an infinite outermost trace, violating fairness. This
establishes the third assertion and completes the proof.
\end{proof}
\end{annexitem}

\begin{annexitem}{annex:p12-klop1987-strategy-results}{thm:p12-klop1987-strategy-results}{Parallel outermost and leftmost normalization with their exact extra hypotheses}
\begin{proof}
\paragraph{Parallel outermost.}
Sequentialize each parallel round by contracting its finite disjoint
set of maximal redexes. A residual trace that remains outermost forever
would reach a round boundary and would be selected and contracted in
the following round. It cannot be infinite. Thus this sequentialized
run is outermost-fair. If the strategy had infinitely many rounds, each
positive round gives an infinite ordinary run. By
Lemma~\ref{thm:p12-klop1987-outermost-projection} it has infinitely
many outermost contractions.

Suppose its input has a finite reduction
\(t=t^{(0)}\to\cdots\to t^{(m)}=n\) to a normal form.
Project the supposed infinite outermost-fair run through the first step
of this finite reduction. The lemma preserves both fairness and
infinitude, so its projection again has infinitely many outermost
contractions. Repeat for each of the finitely many steps of the
normalizing reduction. The last projection would be an infinite
reduction starting at \(n\), which is impossible. Therefore the
parallel-outermost strategy has finitely many rounds on this input;
since a non-normal boundary always admits a positive round, it ends
at the normal form. No NE assumption is used: erasing an inner
projection is allowed, while an outermost contraction cannot disappear
infinitely often among a finite set of disjoint designated residuals.

\paragraph{The exact left-normal position fact.}
In a regular left-normal TRS, a step at a non-leftmost redex preserves
the leftmost redex as the leftmost occurrence. The latter's pattern
persists by regularity: the step is either disjoint to its right or
inside one of its variable instances. An earlier newly created redex
would have to be above that surviving leftmost occurrence; a disjoint
earlier region is unaffected by a step to its right. In the putative
new outer pattern, the surviving old redex must lie in a variable
instance, or non-overlap would fail. All operator occurrences of that
pattern precede this variable leaf by left-normality. The non-leftmost
step lies to the right of the old redex's head, so cannot change any
of those preceding operator occurrences. With a left-linear pattern
there is also no variable-equality test to change. The supposed new
pattern would already have matched before the step, contradicting that
the old occurrence was leftmost. This proves the position fact.

\paragraph{Projection of a leftmost run.}
Let an infinite leftmost run be given, and project it through one
arbitrary initial step, tracking the latter's disjoint designated
residual set \(U_i\). If the original leftmost contraction belongs
to \(U_i\), it cancels. Otherwise none of the designated redexes
can be above it, since an ancestor redex has an earlier head and would
contradict leftmostness. Develop \(U_i\) first. Each designated
step is non-leftmost while the original leftmost redex survives, so
the position fact keeps that redex leftmost throughout the development.
Its contraction in the developed term is a genuine leftmost step.
The updated variable-instance substitution and the canonical
development invariant show that its endpoint is exactly the next
projection boundary. This covers disjoint and enclosing-original-step
cases; an original step strictly inside a designated redex is excluded
by leftmostness. Thus every nonempty projected segment is itself one
leftmost step.

This projection is infinite. Otherwise all original contractions after
some finite row would be designated contractions; their infinite tail
would be a development of that row's finite designated set, contrary
to finite-development termination. After omitting empty segments the
projection is consequently an infinite leftmost run. Project successively
through a finite normalizing reduction, exactly as above. One would
obtain an infinite leftmost run from its normal endpoint. This
contradiction proves that leftmost reduction, equivalently the
leftmost-outermost choice in prefix order, normalizes every normalizable
term in a regular left-normal system. The left-normal position fact
is essential: \(F(C,A)\to F(C,B)\) in the source's example creates
a new earlier root redex after a non-leftmost step when left-normality
fails.

\paragraph{The negative and cofinality columns.}
The regular left-normal system \(Kxy\to x,a\to a\), on
\(Kba\), loops under both innermost strategies although the term
has normal form \(b\). Thus both minus rows are genuine failures
even under the extra restriction. The source's
\(F(x,B)\to D,A\to B,C\to C\) system is regular but not
left-normal, and leftmost-outermost loops on \(F(C,A)\), which
has normal form \(D\). This supplies its unrestricted minus.
The CL invariant and period in Theorem~\ref{thm:p12-klop1987-omega}
show noncofinality of parallel-outermost even for a left-normal regular
system. Leftmost-outermost does not merely sequentialize that period:
it can recreate an enclosing \(S\)-redex before reducing the right arm.
Its own run nevertheless fails cofinality on the same input. Write
\(a=I_{II}=SK(II)\), \(w=Saa\), and
\(X_0=w,X_{j+1}=aX_j\). Its successive boundary terms are
\(wX_0,wX_1,wX_2,\ldots\). From \(wX_j\), the selected steps are
\[
 wX_j\to aX_j(aX_j)
 \to KX_j((II)X_j)(aX_j)\to X_j(aX_j)
 \to^*w(aX_j)=wX_{j+1}.
\]
The last segment peels the \(j\) left-spine \(a\)-applications;
each uses its \(S\)-step followed by its \(K\)-step. Every
\((II)Y\) introduced along the way is in that \(K\)-step's
discarded argument. None of the internal \(II\)-occurrences is
selected: the displayed root or left-spine redex always precedes it.
At every intermediate point all \(S\)-headed codes are unchanged
copies of \(SK(II)\) and \(S(SK(II))(SK(II))\); the selected
\(S\)- and \(K\)-steps remove or copy those codes and do not form
\(SKI\). Hence no point contains \(w_I=S(SKI)(SKI)\).
Every reduct of \(\Omega_I\) does contain \(w_I\), by its
proved invariant, so this descendant reaches no leftmost iterate.
This gives the leftmost plus-but-not-cofinal column separately, without
mistaking its growing left-spine run for the parallel period.
The full-substitution double-plus assertion has the separate complete
proof in Theorem~\ref{thm:p12-klop1987-fair-cofinal}. These deductions
give all five rows without identifying normalization with termination
of arbitrary reductions or with cofinality.
\end{proof}
\end{annexitem}

\begin{annexitem}{annex:p12-klop1987-win-normalization}{thm:p12-klop1987-win-normalization}{An innermost step cannot erase an infinite computation below itself}
\begin{proof}
If a term is SN, repeatedly select an innermost redex until none is
present. Every non-normal finite term has such a redex, by choosing one
minimal under proper containment. SN prevents an infinite sequence of
these choices, so this process ends in a normal form and proves WIN.

Conversely, fix an innermost step \(t\to s\), and suppose an infinite
reduction starts at \(t\). Track the contracted redex by an origin
label along that reduction. Its residuals are disjoint and finite at
every finite stage, by the regular residual construction. In every such
residual, all proper subterms are the original proper subterms: steps
above it can copy or erase the entire redex, and disjoint steps leave it
unchanged. No step can occur strictly inside it, because those proper
subterms were normal to begin with. Hence every surviving residual is
still innermost. Created redexes above it do not change this fact.

Complete the labelled residual set at each stage and project the next
infinite-reduction step past this development. A step at a labelled
residual projects to an empty segment. A disjoint step projects to the
same genuine step. A step above some labelled residuals still applies
after their development: all selected redexes occur in its variable
instances, so developing them updates a consistent matching substitution
and contracting the outer rule gives the common result. This is a
genuine projected step, even if that outer rule erases the labelled
instances. There is no fourth, inner case, since residuals are innermost.
These are the full position cases for the projection.

There must be infinitely many unlabelled contractions in the infinite
sequence. If only finitely many occurred, its final infinite tail would
be a development of a fixed finite labelled set, contradicting
Lemma~\ref{thm:p12-klop1987-developments}. Each of the infinitely many
unlabelled steps projects to a nonempty segment; their concatenation
is an infinite reduction from \(s\). Thus an innermost step preserves
the existence of an infinite reduction without an NE hypothesis.

Now suppose \(t\) has a finite innermost reduction to a normal form.
If \(t\) admitted an infinite reduction, apply this preservation fact
successively along that finite innermost reduction. Its normal endpoint
would admit an infinite reduction, a contradiction. Therefore WIN
implies SN locally; quantification gives the system-wide equivalence.
The restriction to innermost steps replaces NE precisely by excluding
the erasing inner-projection case, rather than by assuming that no rule
can erase an argument.
\end{proof}
\end{annexitem}

\begin{annexitem}{annex:p12-klop1987-fair-cofinal}{thm:p12-klop1987-fair-cofinal}{Finite residual forests, eventual common tails, and complete rounds}
\begin{proof}
First let \(\rho:t=t_0\to t_1\to\cdots\) be an infinite fair
reduction, and fix one redex occurrence \(u\) in \(t\).
The forest of its residuals along successive stages is finitely branching:
every right side is a finite term, so each surviving occurrence has
finitely many copies in the next term. The roots are finite in number
(here one). If residuals of \(u\) existed at arbitrarily late stages,
the elementary finitely branching tree argument would give an infinite
branch: choose a root with descendants at unbounded depths, then a child
with that property, and repeat. That branch is an infinite redex trace,
contrary to fairness. Consequently some finite stage has no residual of
\(u\), and no later stage can recreate its origin label. Newly created
redexes have different origins.

Let \(t\to_u s\) be the fixed step, and let \(U_i\) be its
residual set in \(t_i\). Develop \(U_i\) completely to a term
\(s_i\), with \(s_0=s\). Project each \(t_i\to t_{i+1}\)
past this finite development. The position cases are the same as in
the regular parallel diamond: a labelled contraction cancels; a disjoint
step survives; a step above labels uses the substitution in which those
labels have been developed; and a step below a label is copied into
every retained right-side variable occurrence, possibly giving an empty
segment when that variable is erased. The two resulting routes end at
the same canonical-development endpoint by
Lemma~\ref{thm:p12-klop1987-developments}. Thus
\(s_i\to^*s_{i+1}\) for every \(i\).

At a stage \(N\) where \(U_N\) is empty, \(s_N=t_N\).
For every later stage the projected step is the original step itself,
since there are no labelled residuals left to affect it. Hence the
projection of \(\rho\) through \(t\to_u s\) shares a literal
tail with \(\rho\), including its occurrence choices. It is fair:
an infinite trace in the projected reduction would restrict to an
infinite trace on that common tail, contradicting fairness of \(\rho\).
In particular \(s\to^*t_N\).

For a fixed finite reduction
\(t\to s^{(1)}\to\cdots\to s^{(m)}\), repeat this construction.
The first projection is fair and eventually identical to \(\rho\).
Project that reduction through its next specified step; it again has
a fair common tail with the previous projection. There are only
finitely many specified steps, so, inductively, the final projection
from \(s^{(m)}\) shares a tail with \(\rho\). A term on this
tail is therefore a descendant of \(s^{(m)}\). This proves cofinality
for every finite descendant, without requiring one uniform bound for
all descendants and without assuming NE or termination.

For a finite reduction ending in a normal form, confluence supplies a
join between that endpoint and every other descendant of its input.
Irreducibility forces that join to be the endpoint, proving cofinality
of that completed finite run. The endpoint condition is essential:
finite prefixes satisfy an absence of infinite traces vacuously but need
not yet be cofinal.

Finally sequentialize a full-substitution run into complete development
rounds. A redex trace starting during a round either dies before its
end or reaches the next round's initial term. There it belongs to the
set of all initial redexes, whose complete development ends the trace
before that next round finishes. Thus no infinite redex trace exists.
Each round is finite by the development theorem; a non-normal boundary
has a positive round, even for a rule whose contractum equals its input.
An infinite iteration is therefore an infinite fair ordinary reduction.
A finite completed iteration ends at a normal form. In either case the
preceding argument gives cofinality; every intermediate term of a round
reduces to its boundary endpoint, so cofinality of the sequentialized
run also gives cofinality of the full-substitution iterates themselves.
The already proved cofinal-to-normalizing implication gives normalization.
\end{proof}
\end{annexitem}

\begin{annexitem}{annex:p12-klop1987-omega}{thm:p12-klop1987-omega}{A syntactic invariant for the source's nonreachability assertion}
\begin{proof}
Write \(a=SKI\) and \(w=Saa\), both fixed normal CL terms.
Then \(\Omega_I=ww\). Use four auxiliary constants
\(\mathsf a,\mathsf w,\mathsf i,\mathsf k\), whose interpretations
are \(a,w,I,K\), respectively, and preserve binary application under
interpretation. Define a class of auxiliary values recursively by
\[
 V::=\mathsf w\mid VV\mid\mathsf a V
       \mid\mathsf i V\mid\mathsf k V.
\]
The production \(VV\), applied to \(\mathsf k V\) and another
value, includes \(\mathsf k V V\). Every value contains at least
one occurrence of the auxiliary constant \(\mathsf w\), so its
interpretation contains the literal normal subterm \(w=S(SKI)(SKI)\).

The interpretations of values are closed under CL reduction. To see
this at every position, note first that the fixed codes \(a,w,I,K\)
are normal. A redex crossing a code boundary must be one of these four
forms, with a value \(X\) or values \(X,Y\):
\begin{align*}
 \mathsf w X&\longmapsto\mathsf a X(\mathsf a X),&
 \mathsf a X&\longmapsto\mathsf k X(\mathsf i X),\\
 \mathsf i X&\longmapsto X,&
 \mathsf k X Y&\longmapsto X.
\end{align*}
The first two are exactly the CL \(S\)-rule on the codes
\(Saa\) and \(SKI\); the last two are the \(I\)- and
\(K\)-rules. Proper partial spines of the codes have fewer than the
required arguments and give no further cases. A step inside a value
argument is treated recursively. In particular a code's fixed argument
\(I\) is not itself an applied \(I\)-redex before its enclosing
\(S\)-contraction; that contraction exposes precisely
\(\mathsf i X\), already covered above. The displayed replacements
all yield values. This proves closure by structural induction on the
auxiliary value and exhaustive inspection of its fixed code boundaries.

Since \(\mathsf w\mathsf w\) is a value, every reduct of
\(\Omega_I\) contains the literal subterm \(w\). On the other
hand, \(\Omega_{II}\) has only the codes \(SK(II)\) and
\(S(SK(II))(SK(II))\) beneath its outer application. It contains
no occurrence of \(SKI\), hence no occurrence of \(w\).
Therefore it is not a reduct of \(\Omega_I\).

Reducing the four internal occurrences of \(II\) to \(I\)
in \(\Omega_{II}\) yields \(\Omega_I\). For an arbitrary
parameter \(t\), put \(W_t=SI_tI_t\). At the three successive
parallel-outermost stages, contraction gives
\[
 W_tW_t\to I_tW_t(I_tW_t)
 \to_{\parallel}KW_t(tW_t)(KW_t(tW_t))
 \to_{\parallel}W_tW_t.
\]
At each stage the selected redexes cover the occurrences of the parameter;
no reduction inside \(t\) is selected. This proves the periodic orbit
for \(t=II\). If \(\Omega_I\) reached any point on that orbit,
following the remaining part of the period would reach
\(\Omega_{II}\), contradicting the invariant. Thus this descendant
of the input reaches no strategy iterate. The orbit is not cofinal.
\end{proof}
\end{annexitem}

\begin{annexitem}{annex:p12-klop1987-ne-normalization}{thm:p12-klop1987-ne-normalization}{Projection of an infinite reduction through a non-erasing step}
\begin{proof}
SN implies WN by well-founded induction on reduction: a term with no
successor is normal, and otherwise a successor has a normal-form descendant
by induction. For the converse we prove the stronger preservation fact:
in a regular NE system, if \(t\) admits an infinite reduction and
\(t\to s\), then \(s\) also admits an infinite reduction.

Fix the contracted occurrence \(u\) in \(t\) and give it an origin
label. Along an infinite reduction
\(t=t_0\to t_1\to t_2\to\cdots\), keep its residual set
\(U_i\). Initially \(U_0=\{u\}\). Each residual set is finite
and disjoint. Duplication of an enclosing variable instance copies its
label to disjoint leaves; contraction of a labelled residual removes it;
all other labels survive by the same residual rules as in
Lemma~\ref{thm:p12-klop1987-developments}. Non-overlap prevents an
unlabelled step from destroying a labelled operator pattern. Non-erasure
prevents destruction by erasing the variable instance containing it.
Let \(s_i\) be the endpoint obtained by developing \(U_i\) in
\(t_i\); then \(s_0=s\).

Commute each step \(t_i\to t_{i+1}\) past this development.
If it contracts a labelled residual, the canonical-development endpoint
does not change, so \(s_i=s_{i+1}\). If it is disjoint from the
labelled set, its unchanged occurrence contracts in \(s_i\), giving
\(s_i\to s_{i+1}\). If it is above labelled occurrences, all of
those lie in its variable instances. Developing them first gives a new
matching substitution; contracting the same outer rule then gives exactly
the right side with the selected residuals developed. Thus there is again
a genuine step \(s_i\to s_{i+1}\). If it is below a labelled
occurrence, it lies inside one of that occurrence's variable instances.
Developing the labelled root copies that inner step to every right-side
occurrence of the variable. NE guarantees at least one such occurrence,
and finitely many disjoint copies give a positive finite reduction
\(s_i\to^+s_{i+1}\). Other disjoint labelled roots do not affect
it. These are all relative position cases. Identity-valued rules still
give genuine counted steps; syntactic equality of endpoints does not
make a nonempty sequence empty.

There are infinitely many steps not contracting a labelled residual.
Otherwise, after the last such step, the remaining infinite tail would
be a development of the fixed finite residual set at that stage, contrary
to finite-development termination. Each of the infinitely many unlabelled
steps projects to a positive finite reduction of the \(s_i\), while
labelled steps project to an empty reduction. Concatenating the projected
finite segments therefore gives an infinite reduction from \(s\).
This proves the preservation fact.

If a term with an infinite reduction had a finite reduction to a normal
form, apply the preservation fact at every step of that finite reduction.
Its normal-form endpoint would then admit an infinite reduction, which
is impossible. Therefore WN implies SN locally. Quantifying the local
equivalence over all terms gives the system-wide version. The argument
uses NE precisely to keep the projected inner segment nonempty; without
it the erasing \(K\)-example makes the assertion false.
\end{proof}
\end{annexitem}

\begin{annexitem}{annex:p12-klop1987-ars-core}{thm:p12-klop1987-ars-core}{Core ARS implications with the countability boundary}
\begin{proof}
First suppose CR. Any finite conversion can be joined: induct on its
length, keeping a common descendant of the first endpoint and the current
endpoint. A forward next step and the current joining reduction are a
confluent peak at the current endpoint; join them. A backward next step
can simply be followed by the existing joining reduction. This produces
a join of the conversion endpoints. Conversely the endpoints of a peak
are convertible through its source, so conversion joinability implies CR.
The mixed peak condition implies the same conversion-joining induction;
CR plainly supplies that condition. The conditions on the reflexive
transitive relation in Proposition 1.2 coincide because that relation is
already reflexive and transitive.

If a normal form is convertible with \(b\), conversion joinability
provides a common descendant. The normal form has no step, so that
descendant is the normal form itself; this proves NF. If two normal forms
descend from \(a\), they are convertible; NF makes one reduce to the
other, and irreducibility forces identity. Hence NF implies UN.

Assume UN and WN. For a peak from \(a\) with endpoints \(b,c\),
choose normal forms \(n_b,n_c\) below those endpoints. They are both
below \(a\), so UN identifies them. This proves CR. For any reduction
sequence beginning at \(a\), choose a normal form of \(a\).
Every sequence member has some normal form by WN; it is also a normal
form of \(a\), and UN identifies it with the chosen one. Thus this one
normal form is below every sequence member, proving Ind.

If Ind and Inc held for an infinite sequence \(a_0\to a_1\to\cdots\),
Ind would provide \(b\) below every \(a_i\). Inc would imply
\(|a_i|\leq|b|\) for every \(i\), while the strictly increasing
natural-number sequence \(|a_i|\) is unbounded. This contradiction
proves SN.

For CP implies CR, take a peak from \(a\) ending at \(b,c\).
Cofinality gives sequence members \(a_i,a_j\) below \(b,c\),
respectively. The later member of this same sequence is below both,
giving their join. No countability assumption is used in this direction.
Conversely suppose CR and that \(A\) is countable. Enumerate the
descendants of a fixed \(a\) as \(b_0,b_1,\ldots\), repeating
elements if the set is finite. Begin at \(a\). At stage \(i\),
join the current endpoint with \(b_i\), both descendants of \(a\),
and append the finite reduction from the endpoint to that join. Every
\(b_i\) reaches an endpoint occurring in the concatenated reduction.
If infinitely many actual steps are appended, this is an infinite cofinal
reduction. If only finitely many are appended, its final endpoint remains
a common descendant of every subsequently enumerated element, and the
finite reduction is cofinal. This construction establishes CP for each
\(a\). The enumeration is exactly where countability enters; it is
not supplied for an uncountable ARS by the argument.
\end{proof}
\end{annexitem}

\begin{annexitem}{annex:p12-klop1987-one-rule-puzzle}{thm:p12-klop1987-one-rule-puzzle}{Bottom-up normalization and the self-loop}
\begin{proof}
The term \(q=F(a,F(a,F(b,b)))\) matches the rule at its root with
\(x=a,y=F(b,b)\); its contractum is literally \(q\). Repeating
that genuine rewrite gives an infinite reduction, so the system is not SN.

For WN, normalize finite terms by induction on their tree structure.
Variables and constants have no redex. For \(F(s,t)\), use the
induction hypotheses to reduce its children to normal forms \(s',t'\).
Context closure gives \(F(s,t)\to^*F(s',t')\). If this root is
not a redex, both children are normal, so the entire term is normal.
Otherwise \(s'=a\) and \(t'=F(v,w)\), with normal \(v,w\).
The root contracts to \(F(v,F(v,F(b,b)))\).

If \(v\ne a\) syntactically, none of the new \(F\)-nodes has
left child \(a\), except nodes already inside the normal term \(v\).
There are no redexes there either. Thus this result is normal.
If \(v=a\), the result is the loop term \(q\), but there is a
different finite reduction: contract its proper subterm
\(F(a,F(b,b))\) to \(F(b,F(b,F(b,b)))\), then contract the
remaining outer root. Its final result is
\(F(b,F(b,F(b,b)))\), which is normal because all its \(F\)-nodes
have left child \(b\), not \(a\).
This proves existence of a normal form in every induction case.
It does not prohibit the root self-loop or assert that every reduction
uses this bottom-up choice.
\end{proof}
\end{annexitem}

\begin{annexitem}{annex:p12-klop1987-semantic-un}{thm:p12-klop1987-semantic-un}{From equational soundness to syntactic separation}
\begin{proof}
Evaluation of a rule instance under any assignment gives equal values,
since evaluating a substituted variable defines another assignment to
the variables of the original rule. Equal values remain equal after
evaluation of an outer context, by induction on that context. Thus every
rewrite step is sound in the algebra, and equality of evaluations is
preserved along every finite reduction. If \(t\) reduces to normal
forms \(n,m\), the algebra satisfies \(t=n\) and \(t=m\)
under every assignment, hence satisfies \(n=m\). The stipulated
separation of syntactically distinct normal forms gives \(n=m\)
as terms. This is UN and uses neither WN nor SN nor groundness.
\end{proof}
\end{annexitem}

\begin{annexitem}{annex:p12-klop1987-regular-confluence}{thm:p12-klop1987-regular-confluence}{All positions in the parallel diamond}
\begin{proof}
Represent each parallel contraction by its finite set of disjoint
positions. Prove the diamond by induction on the finite source term.
If neither set includes the root, both sets split into disjoint sets
inside the immediate children. Apply the induction assertion to the
two child contractions in each child, including empty contractions.
Assemble the common child reducts under the unchanged root operator.
The union of the child joining sets is disjoint, and hence gives a
parallel join on each side. For a variable or constant with no selected
root this is the empty contraction case.

If both sets include the root, disjointness says each consists of that
root alone. Non-ambiguity excludes two distinct rule patterns matching
there. A left-linear pattern also determines its matching assignment
uniquely from the concrete input at its variable leaves. Consequently
the two contractions have the same result, and join by empty sets.

Suppose exactly the first set includes the root. It is a singleton root
contraction \(l^\sigma\to r^\sigma\). Every position in the other
set must lie in an instance of a variable of \(l\): otherwise an inner
left-hand-side pattern would overlap a nonvariable position of \(l\),
contrary to regularity. Since \(l\) is left-linear, each variable has
one source occurrence, so the selected inner contractions define a
single substitution \(\sigma'\) with
\(\sigma(x)\to_{\parallel}\sigma'(x)\) for each variable
\(x\) of \(l\). Unaffected variables have empty contractions.
The inner-first endpoint is exactly \(l^{\sigma'}\), not an
inconsistent assignment to repeated variable leaves. Contract its root
to \(r^{\sigma'}\).

The root-first endpoint \(r^\sigma\) reaches \(r^{\sigma'}\)
by copying the selected inner set into every occurrence of its variable
on the right. Distinct variable leaves of the finite right side are
disjoint, and within each copy the selected positions are disjoint.
Thus their union is a disjoint parallel contraction. Duplication gives
more copies; erasure gives no copy and an empty arm. Both are legitimate.
This proves the remaining case; exchanging the two sets proves its
symmetric form. These cases exhaust root and nonroot positions.

The diamond property gives confluence of the reflexive transitive
parallel relation by the finite grid argument: commute the first step
of one path successively past every step of the other, and repeat for
the remaining steps of the first finite path. Each cell is a parallel
diamond; the far corner joins the two endpoints. Each ordinary step is
a singleton parallel step, and each parallel step sequentializes into
ordinary steps. The two reflexive transitive relations therefore agree.
Their confluence proves ordinary CR without termination.
\end{proof}
\end{annexitem}

\begin{annexitem}{annex:p12-klop1987-cofinal-normalizing}{thm:p12-klop1987-cofinal-normalizing}{Cofinality and the role of irreducibility}
\begin{proof}
Let \(n\) be a normal form below \(t\). Cofinality of the strategy
orbit gives an iterate \(\mathbb F^k(t)\) below \(n\).
A normal form admits only the empty finite reduction, so that iterate
is literally \(n\). Thus the strategy normalizes every normalizable
input. If CR holds, two normal forms of \(t\) have a common descendant;
irreducibility identifies each with that descendant, giving uniqueness.
The definitions impose no obligation on merely normalizing strategies
to overtake a reduct of an input that has no normal form. The source's
\(\Omega\) example exploits exactly this additional cofinality demand.
\end{proof}
\end{annexitem}

\begin{annexitem}{annex:p12-klop1987-developments}{thm:p12-klop1987-developments}{Residual labels, termination, and the canonical development}
\begin{proof}
Fix a finite input term and a finite selected set \(U\) of its redex
occurrences. Give each selected occurrence a distinct origin label.
Across a contraction of a selected redex, unselected positions disjoint
from it survive unchanged. A selected redex below it lies in a variable
instance by non-overlap. Its label is copied at every occurrence of that
variable on the right, or erased if the variable is absent there.
A selected ancestor retains its label because rewriting inside one of
its variable instances preserves its left-linear pattern match. No new
label is attached to a newly created, unselected redex. These rules define
the residual selected set after each development step.

Order the finitely many original labels by their depths in the original
selected containment forest: depth zero for selected occurrences without
a selected ancestor, depth one for their nearest selected descendants,
and so on. An ancestor's origin depth is strictly smaller than a
descendant's. In a current term, an ancestor label can contain only
labels originating strictly below it. Contracting a label of depth
\(d\) decreases the number of depth-\(d\) residuals by one, leaves
all counts at smaller depths unchanged, and can change only counts at
larger depths through duplication or erasure. The finite tuple of
nonnegative counts, ordered lexicographically with shallower depths
first, strictly decreases. Lexicographic order on a fixed finite product
of \(\mathbb N\) is well-founded, by induction on the tuple length.
Hence every development terminates; it cannot be made infinite by
duplicating deeper labels. Each step leaves a finite term and finitely
many residual labels.

Define a canonical outcome recursively on the original marked term.
If the root is not selected, retain its operator and apply the construction
to its children. If the root is selected and has match \(l^\sigma\),
all deeper selected occurrences lie in the separate variable instances.
Recursively develop each such instance to \(\widehat\sigma(x)\),
then output \(r^{\widehat\sigma}\). This recursion is well-founded
on proper source subterms. Multiple occurrences on the right use the
same developed variable instance; erasure omits it.

The same canonical operation can be applied to any current labelled term.
It is invariant under a development step. For a step at a selected root,
developing that root first produces exactly the right side with the
residual labels copied or erased as specified; developing its variable
instances afterwards gives \(r^{\widehat\sigma}\), which is also
the recursive definition before that step. For a step below a selected
root, the step occurs within a variable instance. Induction on that
proper subterm leaves its canonical developed value unchanged, and
therefore leaves the enclosing right-side substitution unchanged.
For a step below an unselected root, apply the same induction to its
affected child and retain the other children. These cases establish
invariance for every labelled contraction, including arbitrary outer
contexts.

A complete development has no selected residuals, so its canonical
outcome is its current unlabelled term. Invariance therefore identifies
all complete-development endpoints with the canonical original outcome.
Termination guarantees that repeatedly choosing a residual reaches such
an endpoint. A zero-redex term is fixed. A non-normal finite term has
at least one selected redex when all its redexes are selected, and the
construction contracts at least one actual rule instance, giving a
positive ordinary reduction to the endpoint even if the endpoint happens
to be syntactically the input itself. For an effective rule presentation,
finite matching and the specified recursive construction compute that
endpoint. This establishes well-defined full substitution, but not yet
its cofinality on arbitrary newly created redexes.
\end{proof}
\end{annexitem}
